% !TeX root = ../main.tex
\section{Control flow}

The operational semantics section gave you the full machinery for the structured core
of an imperative language: assignment, sequencing, \texttt{if}, \texttt{while}, and
\texttt{for} (desugared into \texttt{while}). That's enough to run any SIMP program, but
it says nothing about \emph{evaluation order} inside a single expression, about operators
that skip evaluating one of their operands, or about control that jumps somewhere other
than ``the next instruction.'' This section reuses the same $\steppair{}{}{}{}$ /
$\bigstep{}{}{}$ toolbox to cover those cases, then tours the big non-standard
control-transfer mechanisms you'll meet in real languages: exceptions, continuations, and
coroutines. It closes with the historical argument for why control flow is restricted at
all, rather than left as raw \texttt{goto}.

\subsection{sequencing and evaluation order}

Sequencing $\seq{C_1}{C_2}$ composes commands left to right, nothing new there. The
subtlety lives \emph{inside} an expression: when an expression contains more than one
sub-expression with a side effect, does the language guarantee an order for those effects?

Many languages do. Java, JavaScript, and Scheme guarantee strict left-to-right evaluation
of operands and function arguments. C (and, for many constructs, C++) famously does
\emph{not}: the standard leaves the order in which function arguments are evaluated
unspecified, and mixing a read and a write to the same variable without an intervening
\emph{sequence point} is undefined behavior, not just ``order not guaranteed.'' This is a
real-world gotcha, not a theoretical curiosity:

\begin{lstlisting}[language=C]
int x = 0;
int f(void) { return x++; }
int g(void) { return x++; }

int main(void) {
    printf("%d %d\n", f(), g());
    return 0;
}
\end{lstlisting}

One compiler may evaluate \texttt{g()} before \texttt{f()} (printing \texttt{1 0}), another
may go left to right (printing \texttt{0 1}) --- both are conforming C. The classic
textbook example \texttt{i = i++ + ++i;} is worse still: it's not merely
implementation-defined, it's \emph{undefined behavior}, so the compiler is free to do
anything at all, including something that looks nonsensical under any evaluation order.

Formally: SIMP's small-step rule for $E_1 + E_2$ (as you'd write it) picks a fixed order to
reduce $E_1$ and $E_2$, say left operand first. C's semantics for argument evaluation is
instead a rule schema with no fixed order, a nondeterministic choice among orderings that
each compiler resolves once and for all. Nondeterminism in a semantics is not automatically
a bug (the short-circuit rules below are nondeterministic in exactly this controlled,
observable way) but \emph{unspecified} order combined with mutable shared state is exactly
what makes a program's meaning depend on which compiler you happen to be using.

\subsection{short-circuit boolean operators}

Extend the boolean grammar with two operators whose operands may themselves contain
expressions with effects (a function call that mutates the store, say):
\[
  B ::= \ldots \mid B_1 \,\&\&\, B_2 \mid B_1 \,||\, B_2
\]
How much of $B_1 \,\&\&\, B_2$ actually executes is an observable, semantically meaningful
choice, not just an optimization; it determines whether $B_2$'s side effects happen at
all. The small-step rules:

\[
\infer[\textsc{And-False}]
  {\steppair{B_1 \,\&\&\, B_2}{\sigma}{\mathbf{false}}{\sigma'}}
  {\steppair{B_1}{\sigma}{\mathbf{false}}{\sigma'}}
\qquad
\infer[\textsc{And-True}]
  {\steppair{B_1 \,\&\&\, B_2}{\sigma}{B_2}{\sigma'}}
  {\steppair{B_1}{\sigma}{\mathbf{true}}{\sigma'}}
\]

\[
\infer[\textsc{Or-True}]
  {\steppair{B_1 \,||\, B_2}{\sigma}{\mathbf{true}}{\sigma'}}
  {\steppair{B_1}{\sigma}{\mathbf{true}}{\sigma'}}
\qquad
\infer[\textsc{Or-False}]
  {\steppair{B_1 \,||\, B_2}{\sigma}{B_2}{\sigma'}}
  {\steppair{B_1}{\sigma}{\mathbf{false}}{\sigma'}}
\]

\textsc{And-False} and \textsc{Or-True} are the rules that actually short-circuit: they
discard $B_2$ entirely and never reduce it, even if $B_2$ would otherwise get stuck or
diverge. Contrast this with a hypothetical eager, ``bitwise'' \texttt{\&}/\texttt{|} that
always reduces both sides first --- that operator would need only one rule each, but would
get stuck (or loop) whenever the discarded operand would have.

\begin{example}
Take $\deref{x} \neq 0 \,\&\&\, 10 / \deref{x} > 1$ with $\sigma(x) = 0$. The full
derivation:
\[
\infer[\textsc{And-False}]
  {\steppair{\deref{x} \neq 0 \,\&\&\, 10/\deref{x} > 1}{\sigma}{\mathbf{false}}{\sigma}}
  {
    \infer[\textsc{Neq}]
      {\steppair{\deref{x} \neq 0}{\sigma}{\mathbf{false}}{\sigma}}
      {\sigma(x) = 0}
  }
\]
The second operand, $10/\deref{x} > 1$, which would get stuck on division by zero, is never
reduced --- \textsc{And-False} fires the moment the first operand is \textbf{false} and
discards $B_2$ outright. This is exactly the ``guard \&\& use'' idiom you rely on in C,
Java, or Python whenever you write \texttt{x != 0 \&\& 10/x > 1}: it's only safe because
the language's semantics \emph{guarantees} short-circuiting, not because it happens to be
how your compiler behaves today.
\end{example}

\subsection{exceptions: non-local transfer of control}

So far, running a command from state $\sigma$ either yields a final store $\sigma'$
($\bigstep{C}{\sigma}{\sigma'}$) or never yields anything at all (divergence: no rule ever
fires again, exactly the $\bot$ from denotational semantics). Exceptions add a
\emph{third} outcome: $C$ can abort, raising a value, with control jumping straight to the
nearest handler that catches it.

Write $\bigstep{C}{\sigma}{\mathbf{exc}(v),\,\sigma'}$ for ``$C$ aborts, raising value $v$,
having reached store $\sigma'$ at the point of the throw.'' Extend the grammar:
\[
  C ::= \ldots \mid \mathbf{throw}\ E \mid \mathbf{try}\ C\ \mathbf{catch}(x)\ C
\]

\[
\infer[\textsc{Throw}]
  {\bigstep{\mathbf{throw}\ E}{\sigma}{\mathbf{exc}(v),\,\sigma}}
  {\bigstep{E}{\sigma}{v}}
\]

The rule that matters most is how exceptions propagate through sequencing --- an exception
raised by $C_1$ preempts $C_2$ entirely:
\[
\infer[\textsc{Seq-Exc}]
  {\bigstep{\seq{C_1}{C_2}}{\sigma}{\mathbf{exc}(v),\,\sigma'}}
  {\bigstep{C_1}{\sigma}{\mathbf{exc}(v),\,\sigma'}}
\]
Compare this to the ordinary sequencing rule
$\bigstep{C_1}{\sigma}{\sigma'} \ \bigstep{C_2}{\sigma'}{\sigma''}$ over
$\bigstep{\seq{C_1}{C_2}}{\sigma}{\sigma''}$, which only applies when $C_1$ terminates
normally. \textsc{Seq-Exc} and ordinary sequencing are mutually exclusive: $C_1$ can't both
terminate normally \emph{and} raise, so exactly one of the two rules ever applies to a
given $\seq{C_1}{C_2}$.

\texttt{try}/\texttt{catch} has two cases, whether the protected command raises or not:
\[
\infer[\textsc{Try-Catch}]
  {\bigstep{\mathbf{try}\ C\ \mathbf{catch}(x)\ C'}{\sigma}{\sigma''}}
  {\bigstep{C}{\sigma}{\mathbf{exc}(v),\,\sigma'} \\
   \bigstep{C'}{\sigma'[x \mapsto v]}{\sigma''}}
\qquad
\infer[\textsc{Try-Normal}]
  {\bigstep{\mathbf{try}\ C\ \mathbf{catch}(x)\ C'}{\sigma}{\sigma'}}
  {\bigstep{C}{\sigma}{\sigma'}}
\]

\textsc{Try-Catch} \emph{is} stack unwinding, written as an inference rule: control jumps
straight from wherever the exception was thrown, past any pending sequencing, if, or while
bodies in between, to the nearest enclosing handler --- nothing textually between the
\texttt{throw} and the \texttt{catch} ever runs. If $C$ sits nested inside several sequences
and loops, \textsc{Seq-Exc} (plus analogous propagation rules you'd add for \texttt{if} and
\texttt{while} bodies) fires once per level of nesting in the \emph{proof tree}, not once
per step of ordinary execution --- that's the precise sense in which a non-local transfer
``skips'' normal control flow.

\begin{example}
\begin{lstlisting}[language=Python]
def risky(x):
    if x == 0:
        raise ValueError("boom")
    return 10 // x

try:
    r = risky(0)
    print("never printed")
except ValueError as e:
    r = -1
print(r)  # -1
\end{lstlisting}
Tracing it with the rules above: \texttt{risky(0)} raises, so
$\bigstep{\mathbf{throw}\ \text{ValueError}}{\sigma}{\mathbf{exc}(v),\,\sigma}$ by
\textsc{Throw}. The assignment \texttt{r = risky(0)} therefore itself produces an
$\mathbf{exc}$ outcome, and by \textsc{Seq-Exc} the whole sequence
\texttt{r = risky(0); print("never printed")} raises \emph{without ever touching}
\texttt{print}. The enclosing \texttt{try}/\texttt{except} then applies
\textsc{Try-Catch}: \texttt{r = -1} runs in the store where the exception occurred, giving
a final store with $r = -1$, which the last \texttt{print(r)} observes.
\end{example}

\paragraph{Connection to axiomatic semantics.} Before exceptions, a Hoare triple
$\{P\}\, C\, \{Q\}$ implicitly assumed a two-way split: either $C$ terminates in a state
satisfying $Q$, or $C$ diverges (and the triple says nothing at all). Exceptions add a
third outcome that any honest correctness statement about $C$ must account for: $C$ may
(1) terminate normally satisfying $Q$, (2) diverge, or (3) terminate \emph{abnormally} by
raising some exception. A triple that silently ignores case (3) can be misleading ---
$\{x > 0\}\; y := 10 / x\; \{y \cdot x = 10\}$ reads fine as long as you forget that
division can raise. Hoare logics extended for exception-raising languages split the
postcondition in two: an ordinary postcondition $Q$ for normal termination, and one
exceptional postcondition $R_e$ per exception type $e$, giving something like
$\{P\}\, C\, \{Q \mid R_{e_1}, R_{e_2}, \ldots\}$. Composing two commands in sequence then
needs $C_1$'s \emph{normal} postcondition to imply $C_2$'s precondition (exactly the old
sequencing rule), while each of $C_1$'s exceptional postconditions propagates straight past
$C_2$ to whatever encloses the sequence --- a direct mirror of \textsc{Seq-Exc} above, now
at the level of assertions instead of stores.

\subsection{continuations: the rest of the computation, as a value}

Everything above models ``what happens next'' implicitly, in the shape of the derivation
tree: after evaluating $C_1$ in $\seq{C_1}{C_2}$, ``run $C_2$ next'' is baked into the
sequencing rule itself. A \textbf{continuation} makes that implicit ``what happens next''
into an explicit, first-class value --- literally a function representing the rest of the
computation.

Take an ordinary recursive factorial, in Scheme:
\begin{lstlisting}[language=Lisp]
(define (fact n)
  (if (= n 0)
      1
      (* n (fact (- n 1)))))
\end{lstlisting}

Rewriting it in \textbf{continuation-passing style} (CPS) means every function takes one
extra argument, its continuation $k$: a function to \emph{call} with the answer, instead of
returning it.
\begin{lstlisting}[language=Lisp]
(define (fact-cps n k)
  (if (= n 0)
      (k 1)
      (fact-cps (- n 1) (lambda (v) (k (* n v))))))

(fact-cps 5 (lambda (v) v))   ; => 120
\end{lstlisting}
Nothing in \texttt{fact-cps} ever returns a value up the call stack the ordinary way.
Instead of computing \texttt{(* n (fact (- n 1)))} and handing the result back to its
caller, it builds a closure \texttt{(lambda (v) (k (* n v)))} that says ``once you know the
factorial of \texttt{n-1}, multiply it by \texttt{n} and hand \emph{that} to whoever wanted
my answer,'' and passes that closure down as the new continuation for the recursive call.
Unwinding \texttt{(fact-cps 3 (lambda (v) v))} by hand:
\begin{lstlisting}[language=Lisp]
(fact-cps 3 (lambda (v) v))
= (fact-cps 2 (lambda (v) ((lambda (v) v) (* 3 v))))     ; k1
= (fact-cps 1 (lambda (v) (k1 (* 2 v))))                  ; k2
= (fact-cps 0 (lambda (v) (k2 (* 1 v))))                  ; k3
= (k3 1) = (k2 1) = (k1 2) = 6
\end{lstlisting}
We won't give continuations a denotational semantics here --- that means rewriting the
entire meaning function $\llbracket \cdot \rrbracket$ to be continuation-passing throughout,
which is out of scope for this course. The idea to take away is just this: ``the rest of
the computation'' is a value you can name, save, pass around, and even call more than once.
That single idea is what underlies Scheme's \texttt{call/cc}, the desugaring of
\texttt{async}/\texttt{await} in JavaScript, generators, and --- looking back at the
previous subsection --- exception handling itself: a handler is, informally, ``the
continuation to jump to.''

\subsection{coroutines and generators}

Contrast plain call/return: the caller calls $f$, $f$ runs to completion (or raises, or
loops), and the caller resumes with a single value. A \textbf{coroutine} or
\textbf{generator} suspends \emph{in the middle} of its body, handing control back to the
caller together with an intermediate value, and can later be resumed exactly where it left
off, with its local state --- its whole stack frame --- intact.

\begin{lstlisting}[language=Python]
def countdown(n):
    while n > 0:
        yield n
        n -= 1

for i in countdown(3):
    print(i)   # 3, 2, 1
\end{lstlisting}

Every \texttt{yield} suspends \texttt{countdown} mid-loop; the local variable \texttt{n} is
not lost between suspensions --- it's exactly as if the function's continuation at the
\texttt{yield} point were captured and handed to the caller as ``the next thing to run.''
Compare that to a plain function returning a list \texttt{[3, 2, 1]} all at once: with a
generator, producer and consumer interleave, which matters when the sequence is infinite or
expensive to produce eagerly. JavaScript's \texttt{async function*} (async generators) and
\texttt{await} generalize the same idea to suspend across asynchronous operations, not just
loop iterations.

\subsection{structured vs.\ unstructured control flow}

None of this is how early languages worked. Fortran and early BASIC control flow was built
entirely out of \texttt{goto} plus labels; loops and conditionals were things a programmer
assembled by hand out of jumps. Dijkstra's 1968 letter, ``Go To Statement Considered
Harmful,'' argued that unrestricted jumps make it effectively impossible to reason locally
about a program's state: at any given label, ``where did control come from, and what's true
here'' depends on the entire set of jumps that might land there, a set that grows
combinatorially with program size. Structured constructs (sequencing, \texttt{if},
\texttt{while}) restrict control flow to a small set of shapes with exactly one entry and
one exit --- which is precisely why the axiomatic semantics from the previous section works
at all: the Hoare rule for \texttt{while} relies on there being exactly one way in and one
way out of the loop. A \texttt{goto} into the middle of a loop body would break that rule
outright.

Modern languages didn't eliminate non-local control flow, though, they domesticated it.
Exceptions, \texttt{break}/\texttt{continue}, early \texttt{return}, and generators are all
restricted, well-behaved escapes: each has a statically knowable set of possible targets
(the nearest enclosing handler, the nearest enclosing loop, the caller), which is exactly
the property arbitrary \texttt{goto} lacks. Continuations, seen this way, are the escape
hatch of last resort: \texttt{call/cc} gives you the full unrestricted power Dijkstra warned
about, plus a value to reason about, which is why almost no mainstream language exposes it
directly --- but nearly every mainstream language builds specific, tamed continuations
(exception handlers, generators, \texttt{async}/\texttt{await}) straight into its runtime.

\subsection{Exercises}

\begin{itemize}
  \item Give the full small-step derivation, using \textsc{Or-True}/\textsc{Or-False}, for
  $\deref{x} = 0 \,||\, 10/\deref{x} > 1$ with $\sigma(x) = 0$. Which operand never gets
  reduced, and why does that matter?

  \item Suppose you add a non-short-circuiting boolean \textbf{xor} to the grammar, one
  that must evaluate \emph{both} operands no matter what. Write its single small-step rule
  and explain in a sentence or two why \textbf{xor} cannot be given the two-rule treatment
  that $\&\&$ and $||$ got above.

  \item Using \textsc{Throw}, \textsc{Seq-Exc}, and \textsc{Try-Catch}, give a full
  derivation tree for
  \[
    \bigstep{\mathbf{try}\ (\seq{\assign{y}{1}}{\mathbf{throw}\ 99})\ \mathbf{catch}(x)\
    \assign{y}{\deref{x}+1}}{\sigma_0}{\sigma_2}
  \]
  where $\sigma_0(y) = 0$. What is $\sigma_2(y)$?

  \item Booleans need an exception outcome too, not just commands: write the missing
  small-step rule for what happens to $\ift{B}{C_1}{C_2}$ when the guard $B$ itself raises
  before ever reaching a truth value.

  \item Rewrite a function \texttt{sum-to(n)} that adds $1 + 2 + \cdots + n$ recursively
  into continuation-passing style (mirror \texttt{fact-cps} above), then hand-unwind
  \texttt{(sum-to-cps 3 (lambda (v) v))} step by step the way \texttt{fact-cps} was unwound
  in the text.

  \item A colleague claims: ``a \texttt{return} inside a \texttt{try} block, when the
  \texttt{try} also has a \texttt{finally}, is just an ordinary function return, no
  non-local transfer involved.'' Give a short Java example showing this claim is false
  (hint: what does a \texttt{return} inside \texttt{finally} do to a \texttt{return}
  already pending from the \texttt{try}?), and explain in one or two sentences why this
  belongs in this section's list of non-local control transfers.
\end{itemize}
